% ═══════════════════════════════════════════════════════════════════════════
\chapter{Conclusiones y trabajo futuro}
\label{chap:conclusiones}
% ═══════════════════════════════════════════════════════════════════════════

Este trabajo propuso y evaluó una intensidad semiparamétrica de llegada de órdenes agresivas para market making en BTC/USDT, $\lambda_i(t,\delta\mid X_t)=\exp\{g_i(X_t)+h_i(t)+m_i(X_t)f_i(\delta)\}$, y la comparó con Avellaneda--Stoikov y con un Hawkes bivariado reforzados con los mismos hallazgos. Las conclusiones son:

\begin{enumerate}
\item \textbf{La intensidad semiparamétrica predice mejor.} Supera a las referencias en las cinco ventanas del \textit{walk-forward}, en el test y en una simulación posterior de 36 días sin reajustes, en ambos lados del libro, con un AUC de 0,81--0,85. Los intervalos de confianza por días excluyen el cero frente a todas las referencias, salvo frente a E3-fkX en el lado ask, donde la diferencia no es significativa en el test ni en la simulación. Se cumplieron todos los criterios preregistrados del \textit{walk-forward} y del test; en la simulación falló el de calibración tras barridos.
\item \textbf{La ganancia proviene sobre todo del estado del mercado.} Un modelo lineal con las mismas covariables recupera el 90\,\%; la no linealidad añade el resto y, frente a las referencias evaluadas, mejoró la discriminación y la calibración tras barridos. El Hawkes recupera el 60--74\,\% solo cuando tiene dos escalas de memoria.
\item \textbf{La forma en la distancia importa.} Una curva monótona modulada por el estado corrige la mala calibración de $e^{-k\delta}$, y las covariables de barrido corrigen la subestimación del rebote en tres de cuatro situaciones.
\item \textbf{La estimación \textit{data-driven} cambia la política de cotización.} La HJB resuelta numéricamente y aplicada con coeficientes congelados, con la curva de distancia interpolada, se sesga según el cociente de intensidades entre lados y según la pendiente de la curva; con E3 opera entre un 31 y un 38\,\% menos que AS y Hawkes con la misma forma de curva. En la aproximación cerrada GLFT, el nivel de la intensidad solo entra en el término de inventario, porque la distancia base $\gamma^{-1}\log(1+\gamma/k_i)$ no depende de $A_i$.
\item \textbf{Ninguna política es rentable neta de las comisiones evaluadas.} Con la regla de cruce estricto, con cualquier intensidad y cualquiera de las tres políticas, el market making pasivo pierde aproximadamente entre 1,1 y 1,3~bps por ejecución antes de comisiones, en el test y en la simulación; con una cola FIFO aproximada el resultado no cambia. La excepción es la HJB con E3 bajo una ejecución optimista al toque, sin cola delante: obtiene una ganancia bruta pequeña en la simulación (de $+0{,}15$~bps por ejecución con E3 tesis a $+0{,}44$~bps con E3 base) y supera a AS y Hawkes, pero no cubre las comisiones de 7,5--10~bps. Bajo cruce estricto, lo que se captura al cotizar más lejos queda compensado por un movimiento posterior mayor del precio, la mayor parte en el primer segundo: es selección adversa en un mercado con spread de un tick. Con cruce estricto, las diferencias entre modelos son pequeñas; la única significativa es que, con la regla miope, E3 pierde 0,08~bps más por ejecución en la simulación.
\end{enumerate}

La contribución del trabajo es, por tanto, predictiva, y su resultado económico es negativo neto de las comisiones evaluadas. Ambos se sustentan con el mismo protocolo: datos propios auditados, especificaciones congeladas antes de evaluar y criterios fijados de antemano. El periodo de test había sido consultado antes con una versión anterior del modelo, y congelar una versión nueva no recupera su independencia.

\paragraph{Trabajo futuro.}
\begin{itemize}[nosep]
\item Modelar conjuntamente la probabilidad de ejecución y el \textit{markout} posterior, que es la cantidad económicamente relevante. En la fase exploratoria, las covariables disponibles explicaron menos del 3\,\% de la variación del \textit{markout} por ejecución, por lo que su valor con datos de 100~ms es dudoso.
\item Repetir la comparación en mercados con spread más ancho o con comisión maker nula o negativa, y en otras plataformas como Coinbase BTC-USD, donde la ventaja predictiva podría monetizarse.
\item Incorporar $X_t$ como variable de estado en la HJB, en lugar de coeficientes congelados.
\item Medir la prioridad real en la cola con datos del libro por mensaje, y estimar el resultado de la HJB con E3 sin depender de los escenarios de cola de este trabajo.
\item Medir la latencia propia en operación real o capturar desde la región del motor de Binance, para estudiar las carreras de cancelación que los datos de 100~ms no permiten simular.
\end{itemize}
